\documentclass[12pt]{article}
\usepackage[utf8]{inputenc}
\usepackage{amsmath, amssymb}
\usepackage{xcolor}
\usepackage{geometry}
\usepackage{hyperref}
\usepackage{fancyhdr}
\usepackage{enumitem}
\usepackage{minted}
\usepackage{booktabs}
\usepackage{tikz}
\usetikzlibrary{shapes, arrows, positioning}

\geometry{margin=1in}
\hypersetup{colorlinks=true, linkcolor=blue, urlcolor=cyan}

\pagestyle{fancy}
\fancyhf{}
\fancyhead[L]{\textbf{\TOPICTITLE}}
\fancyhead[R]{\thepage}

% ------------------------------- 
% Topic Metadata
% ------------------------------- 
\newcommand{\TOPICTITLE}{Function Programming SKills Introduction}
\title{\TOPICTITLE\\\large Study-Ready Notes}
\author{Compiled by Andrew Photinakis}
\date{\today}

\setlength{\headheight}{15pt}

\begin{document}
\maketitle
\tableofcontents
\newpage

% \section{Function Programming SKills Introduction - Haskell}



\section{Structural Recursion}

Structural recursion is a form of recursion where the implementation follows a template defined by the data structure itself. It typically consists of two main parts:

\begin{itemize}
  \item \textbf{Case Analysis:} The function checks which ``choice'' characterizes the data (e.g., is the list empty or does it have elements?).
  \item \textbf{Recursive Call:} The function calls itself on the ``sub-structure'' of the data.
\end{itemize}

\subsection{Example: List Length}

A list of type $a$ is defined as either an empty list \texttt{[]} or a value $x$ prepended to another list $xs$ (written as \texttt{x:xs}). Calculating the length follows this structure:

\begin{itemize}
  \item \textbf{Base Case:} The length of an empty list \texttt{[]} is $0$.
  \item \textbf{Recursive Step:} The length of \texttt{x:xs} is $1 + \mathrm{len}(xs)$.
\end{itemize}

\section{The Fold Functions}

Functional programming allows us to ``reify,'' or turn into a concrete reusable tool, the abstract patterns of structural recursion. This is done through folds.

\subsection{\texttt{foldr} (Right Fold)}

\texttt{foldr} captures the standard structural recursion pattern.

\begin{itemize}
  \item \textbf{Pattern:}
        \[
          g(x:xs) = f \ x \ (g \ xs)
        \]
        with base case
        \[
          g \ [] = z.
        \]
  \item \textbf{Visualization:} It replaces the list constructors with a function $f$ and the empty list with a starting value $z$, effectively processing from the right.
\end{itemize}

\subsection{\texttt{foldl} (Left Fold)}

\texttt{foldl} represents accumulative structural recursion.

\begin{itemize}
  \item \textbf{Pattern:} It uses an accumulator $a$. For an empty list:
        \[
          g \ [] \ a = a.
        \]
  \item \textbf{Processing:} For a list \texttt{x:xs}, it updates the accumulator and moves to the next element:
        \[
          g(x:xs) \ a = g \ xs \ (f \ a \ x).
        \]
\end{itemize}

\section{Types}

Haskell uses a robust type system to ensure safety and clarity. Common basic types include:

\subsection{Atomic Types}
\begin{itemize}
  \item \texttt{Bool} (true/false)
  \item \texttt{Integer} (arbitrary precision)
  \item \texttt{Int} (fixed precision)
  \item \texttt{Double} (floating point)
  \item \texttt{Char} (single characters)
\end{itemize}

\subsection{Collection Types}
\begin{itemize}
  \item \texttt{String} (lists of characters)
  \item \texttt{[a]} (a list of elements of type $a$)
\end{itemize}

\subsection{Structured Types}
\begin{itemize}
  \item \textbf{Tuples:} $(\alpha_1, \dots, \alpha_n)$ for grouping different types.
  \item \texttt{Maybe a}: Used for optional values.
  \item \texttt{Either a b}: Used for values that can be one of two different types.
\end{itemize}

\section{Type Classes and Overloading}

Type classes define a set of functions that can be applied to different types, allowing for ad-hoc polymorphism (overloading).

\begin{center}
  \begin{tabular}{|l|p{6cm}|p{4cm}|}
    \hline
    \textbf{Type Class} & \textbf{Description}                                     & \textbf{Common Operations}                                                   \\
    \hline
    \texttt{Eq}         & Types that support equality testing.                     & \texttt{==}, \texttt{/=}                                                     \\
    \hline
    \texttt{Ord}        & Types with a total ordering (inherits from \texttt{Eq}). & \texttt{<}, \texttt{<=}, \texttt{>}, \texttt{>=}, \texttt{max}, \texttt{min} \\
    \hline
    \texttt{Show}       & Types that can be converted to strings.                  & \texttt{show}                                                                \\
    \hline
    \texttt{Num}        & Basic numeric types.                                     & \texttt{+}, \texttt{-}, \texttt{*}, \texttt{abs}, \texttt{negate}            \\
    \hline
    \texttt{Fractional} & Types supporting division (inherits from \texttt{Num}).  & \texttt{/}, \texttt{recip}                                                   \\
    \hline
    \texttt{Floating}   & Trigonometric and exponential functions.                 & \texttt{exp}, \texttt{log}, \texttt{sqrt}, \texttt{sin}, \texttt{cos}        \\
    \hline
    \texttt{Foldable}   & Data structures that can be folded.                      & \texttt{foldr}, \texttt{foldl}, \texttt{toList}, \texttt{product}            \\
    \hline
  \end{tabular}
\end{center}


\section{Functional Programming Foundations}

Functional programming (FP) is a paradigm that treats computation as the evaluation of mathematical functions, avoiding changing state and mutable data. In Computer Science and Finance, these concepts are particularly powerful for building predictable high-frequency trading systems and complex financial models.

\subsection{1. Pure Functions}

At the heart of FP is the mathematical notion of a function. A function is \emph{pure} if its output is determined solely by its input values, without observable side effects.

\begin{itemize}
  \item \textbf{Predictability:} Like rule-based mathematical expressions, a pure function produces the same result every time it is called with the same arguments.
  \item \textbf{No Side Effects:} Pure functions do not modify global variables or hidden storage locations, avoiding the issues that arise from implicit state changes.
\end{itemize}

\subsection{2. Immutability}

In FP, data does not change once it is created. Instead of modifying an existing structure, a new structure is produced.

\begin{itemize}
  \item \textbf{Mathematical Grounding:} This aligns with the modern view of functions as mappings (sets or infinite tables).
  \item \textbf{Safety:} Immutability makes computation behave more like mathematics, where a value such as $5$ does not later become $6$.
  \item \textbf{Parallelism:} With no shared mutable state, multithreading becomes simpler and safer.
\end{itemize}

\subsection{3. Recursion (Instead of Loops)}

Because FP avoids state changes (such as incrementing a counter $i$ in a \texttt{for} loop), it relies on recursion to perform repetitive tasks.

\begin{itemize}
  \item \textbf{Structural Recursion:} Recursion that follows the template of the data itself.
  \item \textbf{Example:} To compute the length of a list, consider the base case (empty list) and the recursive case (an element prepended to a sub-list).
  \item \textbf{Tail Recursion:} Tail-recursive formulations improve efficiency and can often be proven correct using mathematical induction.
\end{itemize}

\subsection{4. Higher-Order Functions}

A higher-order function is a function that either takes another function as an argument or returns one as a result.

\begin{itemize}
  \item \textbf{Abstraction:} They allow abstraction of common patterns, such as applying a transformation to every element in a list.
  \item \textbf{Reification:} Functions such as \texttt{foldr} (right fold) and \texttt{foldl} (left fold) reify structural recursion patterns into reusable tools.
\end{itemize}

\subsection{5. Function Composition}

Composition combines two or more functions to produce a new function.

\begin{itemize}
  \item \textbf{Combining Forms:} Emphasis is placed on combining functions rather than sequencing commands.
  \item \textbf{Modular Software:} Function application provides a richer means of combination than sequencing, resulting in succinct and maintainable software.
\end{itemize}

\subsection{Key Idea: Transformations vs.\ State Changes}

In imperative programming, the focus is on how to change the state of the computer (e.g., updating variables or moving pointers).

In functional programming, the focus is on what the data is. Computation is expressed as a sequence of transformations:

\begin{enumerate}
  \item Take a data structure.
  \item Apply a transformation (function).
  \item Produce a new data structure.
\end{enumerate}

As a result, functional programs often resemble their mathematical specifications more closely than imperative programs.


\section{Haskell Syntax and Core Language Features}

Haskell’s syntax is concise and closely mirrors mathematical notation. This design promotes safety and correctness by making programs resemble formal specifications rather than step-by-step machine instructions.

\subsection{1. Defining Functions}

Functions are the primary building blocks in Haskell. A basic definition consists of a name, parameters, an equals sign, and an expression.

\textbf{Syntax:}
\begin{verbatim}
functionName arg1 arg2 = expression
\end{verbatim}

\textbf{Example (Factorial):}
\begin{verbatim}
fact n = if n == 0 then 1 else n * fact (n - 1)
\end{verbatim}

\textbf{Application:}
\begin{verbatim}
fact 5
\end{verbatim}

\textbf{Additional Examples:}

\begin{verbatim}
-- Square of a number
square x = x * x

-- Add three numbers
addThree a b c = a + b + c

-- Fibonacci (simple recursive definition)
fib n = if n <= 1 then n else fib (n-1) + fib (n-2)
\end{verbatim}

\subsection{2. Pattern Matching}

Pattern matching defines behavior based on the “shape” of the input data.

\textbf{Factorial with Pattern Matching:}
\begin{verbatim}
fact 0 = 1
fact n = n * fact (n - 1)
\end{verbatim}

\textbf{Lists:}
\begin{verbatim}
-- Length of a list
len []     = 0
len (x:xs) = 1 + len xs

-- Sum of a list
sumList []     = 0
sumList (x:xs) = x + sumList xs
\end{verbatim}

\textbf{Tuples:}
\begin{verbatim}
first (x, _) = x
second (_, y) = y
\end{verbatim}

\subsection{3. Guards}

Guards provide a readable alternative to nested conditionals.

\textbf{Absolute Value:}
\begin{verbatim}
absolute x
  | x < 0     = -x
  | otherwise = x
\end{verbatim}

\textbf{Additional Examples:}

\begin{verbatim}
grade score
  | score >= 90 = "A"
  | score >= 80 = "B"
  | score >= 70 = "C"
  | score >= 60 = "Pass"
  | otherwise   = "Fail"

sign x
  | x > 0     = 1
  | x == 0    = 0
  | otherwise = -1
\end{verbatim}

\subsection{4. \texttt{let} and \texttt{where} Bindings}

Bindings define local names or helper functions.

\textbf{Using \texttt{where}:}
\begin{verbatim}
circleArea r = pi * rSquare
  where
    rSquare = r * r
\end{verbatim}

\begin{verbatim}
sphereVolume r = (4/3) * pi * r^3
  where
    pi = 3.1415926535
\end{verbatim}

\textbf{Using \texttt{let}:}
\begin{verbatim}
cylinderArea r h =
  let sideArea = 2 * pi * r * h
      topArea  = pi * r^2
  in  sideArea + 2 * topArea
\end{verbatim}

\begin{verbatim}
quadratic a b c =
  let disc = sqrt (b^2 - 4*a*c)
  in  ((-b + disc) / (2*a), (-b - disc) / (2*a))
\end{verbatim}

\subsection{5. Indentation Rules (Layout System)}

Haskell uses whitespace to determine structure.

\textbf{Correctly Indented Example:}
\begin{verbatim}
main = do
  putStrLn "What is your name?"
  name <- getLine
  putStrLn ("Hello " ++ name)
\end{verbatim}

Definitions in a \texttt{where} or \texttt{let} block must align vertically. Spaces are preferred over tabs to ensure consistent formatting.

\subsection{6. Basic Operators and Infix Functions}

\textbf{Arithmetic Operators:}
\begin{verbatim}
+  -  *  /  ^  **
\end{verbatim}

\textbf{Boolean Operators:}
\begin{verbatim}
&&  ||  not
\end{verbatim}

\textbf{Comparison Operators:}
\begin{verbatim}
==  /=  <  <=  >  >=
\end{verbatim}

\textbf{Infix vs.\ Prefix:}

\begin{verbatim}
add x y = x + y

-- Prefix usage
add 5 3

-- Infix usage
5 `add` 3
\end{verbatim}

\textbf{Integer Division vs.\ Fractional Division:}

\begin{verbatim}
10 / 4    -- 2.5  (Fractional)
10 `div` 4  -- 2  (Integral)
\end{verbatim}

\subsection{Practice Exercises}

\begin{enumerate}
  \item \textbf{Pattern Matching:} Write a function \texttt{isZero} that returns \texttt{"Yes"} if the input is 0 and \texttt{"No"} otherwise.
  \item \textbf{Guards:} Define a function \texttt{grade} that returns \texttt{"Pass"} if a score is at least 60 and \texttt{"Fail"} otherwise.
  \item \textbf{Bindings:} Compute the volume of a sphere using a \texttt{where} clause to define $\pi$.
  \item \textbf{Layout:} Rewrite misaligned code so that blocks are properly indented.
  \item \textbf{Operators:} Explain the difference between \texttt{/} (Fractional) and \texttt{div} (Integral).
\end{enumerate}


\section{Haskell Type System and Polymorphism}

Haskell’s type system forms a mathematical foundation that promotes robustness and logical soundness. Strong static typing helps ensure safety and correctness, which is especially important in complex computational systems.

\subsection{1. Basic Types}

Haskell provides several core (primitive) types:

\begin{itemize}
  \item \texttt{Int}: Fixed-precision integer (typically 64-bit).
  \item \texttt{Integer}: Arbitrary-precision integer.
  \item \texttt{Float}: Single-precision floating point.
  \item \texttt{Double}: Double-precision floating point.
  \item \texttt{Bool}: Logical values (\texttt{True}, \texttt{False}).
  \item \texttt{Char}: Single Unicode character (e.g., \texttt{'a'}).
  \item \texttt{String}: List of characters (\texttt{[Char]}).
\end{itemize}

\textbf{Examples:}
\begin{verbatim}
x :: Int
x = 42

bigNumber :: Integer
bigNumber = 9007199254740993

flag :: Bool
flag = True

letter :: Char
letter = 'G'

message :: String
message = "Hello"
\end{verbatim}

\subsection{2. Type Annotations}

Type annotations use \texttt{::} and improve clarity and documentation.

\textbf{Examples:}
\begin{verbatim}
piValue :: Double
piValue = 3.14159

add :: Int -> Int -> Int
add x y = x + y

multiply :: Double -> Double -> Double
multiply x y = x * y

isPositive :: Int -> Bool
isPositive n = n > 0
\end{verbatim}

The arrow \texttt{->} separates input types from the return type. Multiple arrows indicate multiple arguments.

\subsection{3. Type Inference}

Haskell automatically infers the most general type consistent with the code.

\textbf{Examples:}
\begin{verbatim}
x = 5 + 5
-- inferred: x :: Num a => a

square n = n * n
-- inferred: square :: Num a => a -> a

isEven n = n `mod` 2 == 0
-- inferred: isEven :: Integral a => a -> Bool
\end{verbatim}

Type inference prevents invalid operations at compile time.

\subsection{4. Parametric Polymorphism}

Parametric polymorphism allows functions to operate uniformly over any type.

\textbf{Examples:}
\begin{verbatim}
lengthList :: [a] -> Int
lengthList []     = 0
lengthList (_:xs) = 1 + lengthList xs

reverseList :: [a] -> [a]
reverseList []     = []
reverseList (x:xs) = reverseList xs ++ [x]
\end{verbatim}

These functions work for lists of any type because they do not depend on the specific element type.

\subsection{5. Type Variables}

Type variables begin with lowercase letters and represent arbitrary types.

\textbf{Identity Function:}
\begin{verbatim}
id :: a -> a
id x = x
\end{verbatim}

\textbf{Tuple Example:}
\begin{verbatim}
first :: (a, b) -> a
first (x, _) = x

second :: (a, b) -> b
second (_, y) = y
\end{verbatim}

\textbf{Constant Function:}
\begin{verbatim}
constValue :: a -> b -> a
constValue x _ = x
\end{verbatim}

The signature \texttt{a -> b -> a} indicates the function takes two arguments of potentially different types and returns the first.

\subsection{6. Polymorphic Standard Functions}

\begin{verbatim}
head :: [a] -> a
head (x:_) = x

map :: (a -> b) -> [a] -> [b]
map _ []     = []
map f (x:xs) = f x : map f xs

filter :: (a -> Bool) -> [a] -> [a]
filter _ [] = []
filter p (x:xs)
  | p x       = x : filter p xs
  | otherwise = filter p xs
\end{verbatim}

\subsection{Practice Exercises}

\begin{enumerate}
  \item \textbf{Type Identification:} Determine appropriate types for:
        \begin{itemize}
          \item \texttt{'G'} $\rightarrow$ \texttt{Char}
          \item \texttt{9007199254740993} $\rightarrow$ \texttt{Integer}
          \item \texttt{True} $\rightarrow$ \texttt{Bool}
          \item \texttt{6.022} $\rightarrow$ \texttt{Double}
        \end{itemize}

  \item \textbf{Type Annotations:} Write the type signature for:
        \begin{verbatim}
multiply :: Double -> Double -> Double
\end{verbatim}

  \item \textbf{Inference:} For
        \begin{verbatim}
isEven n = n `mod` 2 == 0
\end{verbatim}
        the return type is \texttt{Bool}.

  \item \textbf{Polymorphism:} Provide a type signature for:
        \begin{verbatim}
head :: [a] -> a
\end{verbatim}

  \item \textbf{Type Variables:} Interpret:
        \begin{verbatim}
foo :: a -> b -> a
\end{verbatim}
        This means the function takes two arguments of possibly different types and returns a value of the same type as the first argument.
\end{enumerate}



\section{Lists and Recursion in Functional Programming}

In functional programming, lists are the most fundamental data structure. Because mutable state is avoided, lists are processed using recursion and high-level transformations. These patterns naturally generalize to processing streams of financial data, historical price sequences, or event logs.

\subsection{1. List Syntax and Properties}

A list in Haskell is a homogeneous collection of elements (all elements must have the same type).

\textbf{Constructors:}
\begin{itemize}
  \item \texttt{[]} represents the empty list.
  \item \texttt{:} (cons operator) prepends an element to a list.
\end{itemize}

\textbf{Representation:}
\begin{verbatim}
[1,2,3]  ==  1 : (2 : (3 : []))
\end{verbatim}

\textbf{Examples:}
\begin{verbatim}
numbers :: [Int]
numbers = [10,20,30]

chars :: [Char]
chars = ['a','b','c']

word :: String
word = "Haskell"     -- equivalent to ['H','a','s','k','e','l','l']
\end{verbatim}

\textbf{Common List Operators:}
\begin{verbatim}
head [1,2,3]     -- 1
tail [1,2,3]     -- [2,3]
[1,2] ++ [3,4]   -- [1,2,3,4]
1 : [2,3]        -- [1,2,3]
length [1,2,3]   -- 3
\end{verbatim}

Lists are singly linked structures, so prepending with \texttt{:} is efficient ($O(1)$), while appending with \texttt{++} is linear time.

\subsection{2. Recursively Processing Lists}

Recursive list functions follow the structural recursion template:

\begin{itemize}
  \item Handle the empty case.
  \item Handle the non-empty case \texttt{(x:xs)}.
  \item Recursively process the tail \texttt{xs}.
\end{itemize}

\textbf{Length:}
\begin{verbatim}
len :: [a] -> Int
len []     = 0
len (_:xs) = 1 + len xs
\end{verbatim}

\textbf{Sum:}
\begin{verbatim}
sumList :: Num a => [a] -> a
sumList []     = 0
sumList (x:xs) = x + sumList xs
\end{verbatim}

\textbf{Product:}
\begin{verbatim}
productList :: Num a => [a] -> a
productList []     = 1
productList (x:xs) = x * productList xs
\end{verbatim}

\textbf{Reverse (Naive Version):}
\begin{verbatim}
reverseList :: [a] -> [a]
reverseList []     = []
reverseList (x:xs) = reverseList xs ++ [x]
\end{verbatim}

The naive version is $O(n^2)$ due to repeated use of \texttt{++}.

\textbf{Tail-Recursive Reverse (Efficient):}
\begin{verbatim}
reverseFast :: [a] -> [a]
reverseFast xs = revHelper xs []
  where
    revHelper [] acc     = acc
    revHelper (x:xs) acc = revHelper xs (x:acc)
\end{verbatim}

This version runs in $O(n)$ time.

\subsection{3. The ``Big Four'' Higher-Order Functions}

Higher-order functions abstract recursive patterns.

\subsubsection{map}

Applies a function to each element.

\begin{verbatim}
map :: (a -> b) -> [a] -> [b]

map _ []     = []
map f (x:xs) = f x : map f xs
\end{verbatim}

\textbf{Examples:}
\begin{verbatim}
map (+1) [1,2,3]         -- [2,3,4]
map (*1.1) prices        -- apply 10% increase
map length ["AI","ML"]   -- [2,2]
\end{verbatim}

\subsubsection{filter}

Selects elements satisfying a predicate.

\begin{verbatim}
filter :: (a -> Bool) -> [a] -> [a]
\end{verbatim}

\textbf{Examples:}
\begin{verbatim}
filter even [1..10]          -- [2,4,6,8,10]
filter (>100) prices         -- prices above 100
filter (\x -> x < 0) deltas  -- negative returns
\end{verbatim}

\subsubsection{foldr}

Processes from the right.

\begin{verbatim}
foldr :: (a -> b -> b) -> b -> [a] -> b

foldr f z []     = z
foldr f z (x:xs) = f x (foldr f z xs)
\end{verbatim}

\textbf{Examples:}
\begin{verbatim}
foldr (+) 0 [1,2,3]      -- 6
foldr (:) [] [1,2,3]     -- [1,2,3]
foldr (++) [] [[1],[2]]  -- [1,2]
\end{verbatim}

\subsubsection{foldl}

Processes from the left with an accumulator.

\begin{verbatim}
foldl :: (b -> a -> b) -> b -> [a] -> b

foldl f acc []     = acc
foldl f acc (x:xs) = foldl f (f acc x) xs
\end{verbatim}

\textbf{Examples:}
\begin{verbatim}
foldl (+) 0 [1,2,3]      -- 6
foldl (*) 1 [1,2,3,4]    -- 24
foldl (\acc x -> x:acc) [] [1,2,3]  -- [3,2,1]
\end{verbatim}

\subsection{4. List Comprehensions}

List comprehensions provide declarative list construction similar to set-builder notation.

\textbf{Syntax:}
\begin{verbatim}
[ expression | element <- list, qualifier ]
\end{verbatim}

\textbf{Examples:}

\begin{verbatim}
-- Squares of even numbers
[ x*x | x <- xs, even x ]

-- All pairs (x,y)
[ (x,y) | x <- [1,2], y <- [3,4] ]

-- Filter and transform
[ price * 1.05 | price <- prices, price > 100 ]

-- Cartesian product with condition
[ (x,y) | x <- [1..5], y <- [1..5], x < y ]
\end{verbatim}

\subsection{Key Insight}

List processing in functional programming emphasizes:

\begin{itemize}
  \item Structural recursion over mutation.
  \item Declarative transformations over step-by-step state updates.
  \item Reusable abstractions (\texttt{map}, \texttt{filter}, \texttt{fold}) instead of manual loops.
\end{itemize}

These patterns scale naturally from small academic examples to large-scale data transformation pipelines.



\section{Algebraic Data Types (ADTs) and Data Modeling}

In Haskell, Algebraic Data Types (ADTs) are the primary mechanism for modeling domain concepts. They allow complex structures to be built from simpler ones using:

\begin{itemize}
  \item \textbf{Products (AND)} — combining multiple pieces of data.
  \item \textbf{Sums (OR)} — choosing between alternative forms.
\end{itemize}

ADTs enable precise modeling of real-world systems and make many categories of bugs impossible by construction.

\subsection{1. \texttt{data} Declarations}

New types are introduced using the \texttt{data} keyword.

\textbf{Syntax:}
\begin{verbatim}
data TypeName = ConstructorName Type1 Type2 ...
\end{verbatim}

\textbf{Rules:}
\begin{itemize}
  \item Type names begin with an uppercase letter.
  \item Constructor names also begin with an uppercase letter.
\end{itemize}

\textbf{Simple Example:}
\begin{verbatim}
data Color = Red | Green | Blue
\end{verbatim}

Here, \texttt{Color} is a type with three possible values.

\subsection{2. Product Types (AND)}

A product type combines multiple values into one structure.

\textbf{Example: 2D Point}
\begin{verbatim}
data Point = Point Double Double
\end{verbatim}

This defines a type \texttt{Point} consisting of two \texttt{Double} values.

\textbf{Usage:}
\begin{verbatim}
myPoint :: Point
myPoint = Point 10.5 20.0
\end{verbatim}

\textbf{Pattern Matching:}
\begin{verbatim}
distanceFromOrigin :: Point -> Double
distanceFromOrigin (Point x y) = sqrt (x^2 + y^2)
\end{verbatim}

\textbf{Another Product Example:}
\begin{verbatim}
data Stock = Stock String Double
\end{verbatim}

\begin{verbatim}
apple = Stock "AAPL" 195.25
\end{verbatim}

\subsection{3. Sum Types (OR)}

A sum type represents a value that can take one of several forms.

\textbf{Example: Order Types}
\begin{verbatim}
data OrderType
  = Market
  | Limit Double
  | StopLoss Double
\end{verbatim}

An \texttt{OrderType} is either:
\begin{itemize}
  \item \texttt{Market}
  \item \texttt{Limit price}
  \item \texttt{StopLoss price}
\end{itemize}

\textbf{Another Sum Example:}
\begin{verbatim}
data TradeSignal = Buy | Sell | Hold
\end{verbatim}

\subsection{4. Record Syntax}

Record syntax names fields and automatically creates accessor functions.

\textbf{Example: Trader Profile}
\begin{verbatim}
data Trader = Trader
  { name    :: String
  , balance :: Double
  , active  :: Bool
  }
\end{verbatim}

\textbf{Usage:}
\begin{verbatim}
alice :: Trader
alice = Trader "Alice" 100000.0 True

currentBalance = balance alice
\end{verbatim}

\textbf{Record Example: Trade}
\begin{verbatim}
data Trade = Trade
  { ticker   :: String
  , quantity :: Int
  , price    :: Double
  }
\end{verbatim}

\subsection{5. Pattern Matching on Custom Types}

Pattern matching extracts values from constructors.

\textbf{Processing Orders:}
\begin{verbatim}
describeOrder :: OrderType -> String
describeOrder Market =
  "Executing at market price."
describeOrder (Limit price) =
  "Waiting for price " ++ show price
describeOrder (StopLoss p) =
  "Stop loss set at " ++ show p
\end{verbatim}

\textbf{Evaluating a Trade:}
\begin{verbatim}
evaluateTrade :: Trade -> Double
evaluateTrade (Trade _ qty p) =
  fromIntegral qty * p
\end{verbatim}

\textbf{Pattern Matching with Records:}
\begin{verbatim}
evaluateTrade' :: Trade -> Double
evaluateTrade' trade =
  price trade * fromIntegral (quantity trade)
\end{verbatim}

\subsection{6. Parametric ADTs}

ADTs can also be polymorphic.

\textbf{Example: Maybe}
\begin{verbatim}
data Maybe a = Nothing | Just a
\end{verbatim}

\textbf{Safe Price Lookup:}
\begin{verbatim}
safePrice :: Bool -> Double -> Maybe Double
safePrice available p =
  if available then Just p else Nothing
\end{verbatim}

\textbf{Handling Maybe:}
\begin{verbatim}
displayPrice :: Maybe Double -> String
displayPrice Nothing  = "Price unavailable"
displayPrice (Just p) = "Price: " ++ show p
\end{verbatim}

\subsection{7. Making Impossible States Unrepresentable}

ADTs improve safety by forcing all cases to be handled.

\textbf{Example:}
\begin{verbatim}
data Transaction
  = Deposit Double
  | Withdrawal Double
  | CheckBalance
\end{verbatim}

Because the compiler requires exhaustive pattern matching, every transaction type must be handled:

\begin{verbatim}
processTransaction :: Transaction -> String
processTransaction (Deposit amt) =
  "Depositing " ++ show amt
processTransaction (Withdrawal amt) =
  "Withdrawing " ++ show amt
processTransaction CheckBalance =
  "Checking balance"
\end{verbatim}

There is no possibility of an invalid transaction shape.

\subsection{8. Design Tradeoffs}

\begin{itemize}
  \item In object-oriented systems, adding new data variants is easy, but adding new operations often requires modifying many classes.
  \item In functional systems, adding new functions over existing data is easy, but adding new constructors requires updating pattern matches.
\end{itemize}

\subsection{Practice Exercises}

\begin{enumerate}
  \item Define:
        \begin{verbatim}
data Transaction
  = Deposit Double
  | Withdrawal Double
  | CheckBalance
\end{verbatim}

  \item Define:
        \begin{verbatim}
data Stock = Stock String Double
\end{verbatim}

  \item Define:
        \begin{verbatim}
data Trade = Trade
  { ticker   :: String
  , quantity :: Int
  , price    :: Double
  }
\end{verbatim}

  \item Write:
        \begin{verbatim}
evaluateTrade :: Trade -> Double
\end{verbatim}

  \item Explain why:
        \begin{verbatim}
Maybe Double
\end{verbatim}
        is safer than using a null value for a missing stock price.
\end{enumerate}

\textbf{Answer (Conceptual):} \texttt{Maybe} forces explicit handling of absence (\texttt{Nothing}) versus presence (\texttt{Just value}). This prevents null-reference errors and ensures all cases are checked at compile time.


\section{Typeclasses and Ad-Hoc Polymorphism}

In Haskell, typeclasses define a set of behaviors (functions) that different types can implement. They enable \emph{ad-hoc polymorphism}, meaning the same function name can operate differently depending on the type of its arguments.

A useful way to think about typeclasses is as formal interfaces or regulations: they specify \emph{what} a type must be able to do, but not \emph{how} it does it.

\subsection{1. Core Standard Typeclasses}

Haskell provides several fundamental typeclasses used throughout the language.

\subsubsection{Eq (Equality)}

Defines equality comparison.

\begin{verbatim}
class Eq a where
  (==) :: a -> a -> Bool
  (/=) :: a -> a -> Bool
\end{verbatim}

\textbf{Examples:}
\begin{verbatim}
5 == 5        -- True
'a' /= 'b'    -- True
True == False -- False
\end{verbatim}

If a type is an instance of \texttt{Eq}, its values can be compared for equality.

\subsubsection{Ord (Ordering)}

Defines total ordering and inherits from \texttt{Eq}.

\begin{verbatim}
class Eq a => Ord a where
  (<), (<=), (>), (>=) :: a -> a -> Bool
  compare :: a -> a -> Ordering
  max, min :: a -> a -> a
\end{verbatim}

\textbf{Examples:}
\begin{verbatim}
'a' < 'b'      -- True
max 10 20      -- 20
compare 5 3    -- GT
\end{verbatim}

Any type used in sorting must be an instance of \texttt{Ord}.

\subsubsection{Show and Read}

These convert between values and strings.

\begin{verbatim}
show 123         -- "123"
read "123" :: Int -- 123
\end{verbatim}

Most basic types are instances of both.

\subsubsection{Num (Numeric)}

Provides basic arithmetic.

\begin{verbatim}
class Num a where
  (+), (-), (*) :: a -> a -> a
  negate :: a -> a
  abs, signum :: a -> a
  fromInteger :: Integer -> a
\end{verbatim}

\textbf{Hierarchy:}
\begin{itemize}
  \item \texttt{Fractional} extends \texttt{Num} (adds \texttt{/}, \texttt{recip})
  \item \texttt{Floating} extends \texttt{Fractional} (adds \texttt{sin}, \texttt{log}, etc.)
\end{itemize}

\textbf{Example:}
\begin{verbatim}
10 + 5      -- works (Num)
10 / 5      -- requires Fractional
\end{verbatim}

\subsection{2. Typeclass Constraints in Signatures}

Type constraints appear before the \texttt{=>} symbol.

\begin{verbatim}
sum :: (Foldable t, Num a) => t a -> a
\end{verbatim}

This means:
\begin{itemize}
  \item \texttt{t} must be \texttt{Foldable}
  \item \texttt{a} must be \texttt{Num}
\end{itemize}

Another example:

\begin{verbatim}
sort :: Ord a => [a] -> [a]
\end{verbatim}

The list elements must be orderable.

\subsection{3. Writing Your Own Typeclass}

You can define custom behavior using the \texttt{class} keyword.

\textbf{Example:}
\begin{verbatim}
class Priced a where
  getPrice :: a -> Double
\end{verbatim}

This states that any type belonging to \texttt{Priced} must implement \texttt{getPrice}.

\subsection{4. Instance Declarations}

An \texttt{instance} connects a specific type to a typeclass.

\textbf{Example:}
\begin{verbatim}
data Stock = Stock String Double

instance Priced Stock where
  getPrice (Stock _ p) = p
\end{verbatim}

Now any function requiring a \texttt{Priced} type can accept a \texttt{Stock}.

\subsection{5. Making Custom Types Instances of Standard Classes}

You can also make your own types instances of built-in typeclasses.

\textbf{Example:}
\begin{verbatim}
data TradeSignal = Buy | Sell | Hold

instance Show TradeSignal where
  show Buy  = "BUY"
  show Sell = "SELL"
  show Hold = "HOLD"
\end{verbatim}

\textbf{Example with Eq:}
\begin{verbatim}
data Position = Long | Short

instance Eq Position where
  Long  == Long  = True
  Short == Short = True
  _     == _     = False
\end{verbatim}

\subsection{6. Polymorphic Functions with Constraints}

\begin{verbatim}
isEqual :: Eq a => a -> a -> Bool
isEqual x y = x == y
\end{verbatim}

\begin{verbatim}
maximumValue :: Ord a => [a] -> a
maximumValue [x]    = x
maximumValue (x:xs) = max x (maximumValue xs)
\end{verbatim}

\subsection{7. Why Typeclasses Matter}

\begin{itemize}
  \item They enforce compile-time guarantees.
  \item They clarify what operations are valid on a type.
  \item They make error messages understandable when constraints are violated.
\end{itemize}

Example error:

\begin{verbatim}
No instance for (Eq a) arising from a use of ‘==’
\end{verbatim}

This means the type used does not implement \texttt{Eq}.

\subsection{Practice Exercises}

\begin{enumerate}
  \item \textbf{Sorting Constraint:} You need \texttt{Ord a =>} to sort a list.

  \item \textbf{Division Error:} The operator \texttt{/} belongs to \texttt{Fractional}, not \texttt{Num}. Since \texttt{Int} is not \texttt{Fractional}, \texttt{Int / Int} fails.

  \item \textbf{Custom Typeclass:}
        \begin{verbatim}
class Representable a where
  asString :: a -> String
\end{verbatim}

  \item \textbf{Instance for Bool:}
        \begin{verbatim}
instance Representable Bool where
  asString True  = "YES"
  asString False = "NO"
\end{verbatim}

  \item \textbf{Equality Function:}
        \begin{verbatim}
isEqual :: Eq a => a -> a -> Bool
\end{verbatim}
        The required constraint is \texttt{Eq a}.
\end{enumerate}

Typeclasses form the backbone of Haskell’s abstraction system, enabling flexible yet strongly typed polymorphism.

% template
\begin{itemize}
  \item X
        \begin{enumerate}
          \item Y
        \end{enumerate}
\end{itemize}

\end{document}